%%%%%%%%%%%%%%%%%%%%%%%%%%%%% Define Article %%%%%%%%%%%%%%%%%%%%%%%%%%%%%%%%%%
\documentclass{article}
%%%%%%%%%%%%%%%%%%%%%%%%%%%%%%%%%%%%%%%%%%%%%%%%%%%%%%%%%%%%%%%%%%%%%%%%%%%%%%%

%%%%%%%%%%%%%%%%%%%%%%%%%%%%% Using Packages %%%%%%%%%%%%%%%%%%%%%%%%%%%%%%%%%%
\usepackage{geometry}
\usepackage{graphicx}
\usepackage{amssymb}
\usepackage{amsmath}
\usepackage{amsthm}
\usepackage{empheq}
\usepackage{mdframed}
\usepackage{booktabs}
\usepackage{lipsum}
\usepackage{graphicx}
\usepackage{color}
\usepackage{psfrag}
\usepackage{pgfplots}
\usepackage{bm}
%%%%%%%%%%%%%%%%%%%%%%%%%%%%%%%%%%%%%%%%%%%%%%%%%%%%%%%%%%%%%%%%%%%%%%%%%%%%%%%

% Other Settings

%%%%%%%%%%%%%%%%%%%%%%%%%% Page Setting %%%%%%%%%%%%%%%%%%%%%%%%%%%%%%%%%%%%%%%
\geometry{a4paper}

%%%%%%%%%%%%%%%%%%%%%%%%%% Define some useful colors %%%%%%%%%%%%%%%%%%%%%%%%%%
\definecolor{ocre}{RGB}{243,102,25}
\definecolor{mygray}{RGB}{243,243,244}
\definecolor{deepGreen}{RGB}{26,111,0}
\definecolor{shallowGreen}{RGB}{235,255,255}
\definecolor{deepBlue}{RGB}{61,124,222}
\definecolor{shallowBlue}{RGB}{235,249,255}
%%%%%%%%%%%%%%%%%%%%%%%%%%%%%%%%%%%%%%%%%%%%%%%%%%%%%%%%%%%%%%%%%%%%%%%%%%%%%%%

%%%%%%%%%%%%%%%%%%%%%%%%%% Define an orangebox command %%%%%%%%%%%%%%%%%%%%%%%%
\newcommand\orangebox[1]{\fcolorbox{ocre}{mygray}{\hspace{1em}#1\hspace{1em}}}
%%%%%%%%%%%%%%%%%%%%%%%%%%%%%%%%%%%%%%%%%%%%%%%%%%%%%%%%%%%%%%%%%%%%%%%%%%%%%%%

%%%%%%%%%%%%%%%%%%%%%%%%%%%% English Environments %%%%%%%%%%%%%%%%%%%%%%%%%%%%%
\newtheoremstyle{mytheoremstyle}{3pt}{3pt}{\normalfont}{0cm}{\rmfamily\bfseries}{}{1em}{{\color{black}\thmname{#1}~\thmnumber{#2}}\thmnote{\,--\,#3}}
\newtheoremstyle{myproblemstyle}{3pt}{3pt}{\normalfont}{0cm}{\rmfamily\bfseries}{}{1em}{{\color{black}\thmname{#1}~\thmnumber{#2}}\thmnote{\,--\,#3}}
\theoremstyle{mytheoremstyle}
\newmdtheoremenv[linewidth=1pt,backgroundcolor=shallowGreen,linecolor=deepGreen,leftmargin=0pt,innerleftmargin=20pt,innerrightmargin=20pt,]{theorem}{Theorem}[section]
\theoremstyle{mytheoremstyle}
\newmdtheoremenv[linewidth=1pt,backgroundcolor=shallowBlue,linecolor=deepBlue,leftmargin=0pt,innerleftmargin=20pt,innerrightmargin=20pt,]{definition}{Definition}[section]
\theoremstyle{myproblemstyle}
\newmdtheoremenv[linecolor=black,leftmargin=0pt,innerleftmargin=10pt,innerrightmargin=10pt,]{problem}{Problem}[section]
%%%%%%%%%%%%%%%%%%%%%%%%%%%%%%%%%%%%%%%%%%%%%%%%%%%%%%%%%%%%%%%%%%%%%%%%%%%%%%%

%%%%%%%%%%%%%%%%%%%%%%%%%%%%%%% Plotting Settings %%%%%%%%%%%%%%%%%%%%%%%%%%%%%
\usepgfplotslibrary{colorbrewer}
\pgfplotsset{width=8cm,compat=1.9}
%%%%%%%%%%%%%%%%%%%%%%%%%%%%%%%%%%%%%%%%%%%%%%%%%%%%%%%%%%%%%%%%%%%%%%%%%%%%%%%

%%%%%%%%%%%%%%%%%%%%%%%%%%%%%%% Title & Author %%%%%%%%%%%%%%%%%%%%%%%%%%%%%%%%
\title{Pythagorean Triplet}
\author{Spencer Veatch}
%%%%%%%%%%%%%%%%%%%%%%%%%%%%%%%%%%%%%%%%%%%%%%%%%%%%%%%%%%%%%%%%%%%%%%%%%%%%%%%

\begin{document}
    \maketitle
    \begin{problem}[Special Pythagorean Triplet]

       A Pythagorean Triplet is a set of three natural numbers \( a < b < c  \), for which: \[
        a^2 + b^2 = c^2
       \] 
       For example, \( 3^2 + 4^2 = 9 + 16 = 25 = 5^2 \). There exists exactly one Pythagorean Triplet for which \( a + b
       + c = 1000 \). Find the product \( abc \).
    \end{problem}
    \subsection{Reducing The Search Space}\label{sub:Reducing The Search Space} % (fold)
    
    Since this problem gives both a sum constraint and a Pythagorean constraint, we can reduce the search space
    algebraically. Formally, the problem asks us to find the Pythagorean Triplet \( (a,b,c) \) satisfying the system: 
    \[
        \begin{align}
            a^2 + b^2 &= c^2 \\ 
            a + b + c &= 1000 
        \end{align}
    \] 
    Given the ordering constraint \( a < b < c \). Using the sum constraint, we can solve for \( c \): \[
        c = 1000 - a - b
    \]
    which can be substituted into the Pythagorean Equation as \[
        \begin{align*}
            a^2 + b^2 &= (1000 - a -b)^2 \\
            a^2 + b^2 &= (1000000 + a^2 + b^2 -2000a - 2000b + 2ab) \\
            0 &= 1000000 - 2000a - 2000b + 2ab \\
            0 &= 500000 - 1000a - 1000b + ab \\
        \end{align*}
    \] 
    This representation can be rearranged and factored with respect to \( b \) as: 
    \[
        \begin{align}
            ab - 1000a - 1000b &= -500000 \\
            b(a - 1000) &= 1000a - 500000 \\
            b&=  \frac{1000a - 500000}{(a - 1000)} \\
        \end{align}
    \] 
    Thus, for each admissible value of \( a \), the corresponding value of \( b \) can be uniquely determined.
    % subsection Reducing The Search Space (end)
    \subsection{Bounding the Domain}\label{sub:Bounding the Domain} % (fold)
    Given the inequality \[
        a < b < c
    \] 
    and the knowledge that \( a + b + c = 1000 \), it follows that: \[
        \begin{align}
            3a &< 1000 \\ 
            a &< \frac{1000}{3} \\
        \end{align}
    \] 
    Ergo, the search is reduced to the finite interval: \[
        [1 \le a \le 333].
    \] 
    The algorithmic solution must simply iterate over all possible values of \( a \), computing the corresponding values
    of \( b \) and \( c \), the verifying whether \( a^2 + b^2 = c^2 \) holds.
    % subsection Bounding the Domain (end)
\end{document}
